\documentclass[11pt,a4paper]{article}
\usepackage[margin=25mm]{geometry}
\usepackage{amsmath,amssymb,amsthm,booktabs,hyperref}
\newtheorem{lemma}{Lemma}\newtheorem{theorem}[lemma]{Theorem}
\title{O'Nan map types involving factors 5 and 7}
\author{Human Compression project}\date{}
\begin{document}\maketitle
\begin{abstract}
We determine twelve entries of the O'Nan map enumeration from the
previously acquired character counts. A short subgroup argument shows
that every nongenerating pair of these types has exactly one maximal
overgroup. The existing union bound is consequently an equality.
Ten of the twelve entries were not previously determined by the
accepted counting components. No new group enumeration is needed.
\end{abstract}

\section{A subgroup forced by two prime divisors}
Let $G=O'N$. We use the complete maximal-subgroup classification and
the acquired subgroup roster as explicit external premises. Rows in
the latter refer to its fixed 581-class edition. Among all maximal
subgroup classes of $G$, exactly four have order divisible by 35:
\begin{center}\begin{tabular}{rrl}\toprule
row&order&type\\\midrule
6&175560&$J_1$\\
7&161280&$4.L_3(4):2$\\
43,44&2520&$A_7$ (two classes)\\\bottomrule
\end{tabular}\end{center}
Write $M$ for a subgroup in row 7. Its row-8 maximal subgroup $E$
has order 80640, hence is normal of index two. The exact central
extension label is not needed in the following argument.

The same roster gives the orders of all maximal subgroups of the
three groups $J_1$, $A_7$ and $E$:
\begin{center}\begin{tabular}{ll}\toprule
group&maximal-subgroup orders, with class repetitions\\\midrule
$J_1$&$660,168,120,114,110,60,42$\\
$A_7$&$360,168,168,120,72$\\
$E$&$3840,1440,1440,1440,672,672,672,288$\\\bottomrule
\end{tabular}\end{center}
None of these orders is divisible by 35. Thus a subgroup containing
elements of orders 5 and 7 cannot be proper in any of these three
groups: a proper subgroup would lie in one of the listed maximals,
contrary to Lagrange's theorem. These are complete lists, not selected
examples of maximal subgroups.

\begin{lemma}\label{unique}
Every proper subgroup $H<G$ containing elements of orders 5 and 7
has a unique maximal overgroup in $G$.
\end{lemma}
\begin{proof}
Choose a maximal overgroup $K$. Its order is divisible by 35, so it
belongs to one of the four classes above. If $K\cong J_1$ or $A_7$,
the preceding observation forces $H=K$.

Otherwise conjugate so that $K=M$. The quotient $M/E$ has order two.
Every element of odd order in $H$ therefore belongs to $H\cap E$.
In particular $H\cap E$ contains elements of orders 5 and 7.
The maximal-subgroup list of $E$ forces $H\cap E=E$, whence
$E\leq H\leq M$. As $[M:E]=2$, either $H=E$ or $H=M$.
The case $H=M$ is immediate; it remains to prove that $E$ has a
unique maximal overgroup.

Every maximal overgroup $K'$ of $E$ has order divisible by 80640.
The complete maximal-subgroup list of $G$ shows that only the class
of $M$ has this divisibility. Hence $[K':E]=2$, so $K'$ normalizes
$E$. Since $E\lhd M$, the normalizer $N_G(E)$ contains $M$.
It is not $G$, because $E$ is nontrivial and proper and $G$ is simple.
Maximality of $M$ gives $N_G(E)=M$, and hence $K'\leq M$.
Therefore $K'=M$, as required.
\end{proof}
This last step is the normalizer argument used in the earlier exact
columns. In particular it does not identify two subgroup embeddings
merely because they have equal orders.

\section{The twelve exact values}
For a finite group $L$, put
\[
 X_L(p,q)=\{(r,s)\in L^2:|r|=p,\ |s|=q,\ |rs|=2\},
 \qquad x_L(p,q)=|X_L(p,q)|.
\]
Let $e_G(p,q)$ count the pairs generating $G$.
Suppose one of $p,q$ is divisible by 5 and the other by 7. Suitable
powers of $r,s$ then have orders 5 and 7. Every nongenerating pair
in $X_G(p,q)$ satisfies Lemma~\ref{unique}. If $M_j$ represent all
maximal conjugacy classes of $G$, the already established exact
subtraction principle gives
\[
 e_G(p,q)=x_G(p,q)-\sum_j [G:M_j]x_{M_j}(p,q).         \tag{1}
\]
For completeness, each nongenerating pair belongs to a unique
maximal subgroup, so these sets of pairs are disjoint. Each maximal
subgroup is self-normalizing in the simple group $G$, making the
number of conjugates precisely $[G:M_j]$. This proves (1).
Terms for the other nine maximal classes vanish by Lagrange's theorem.

The exact character evaluations on the right side of (1) are already
acquired and independently reviewed. We reuse those integers, including
all conjugacy-class multiplicities. Automorphisms act freely on generating
pairs, so division by $|\operatorname{Aut}(G)|=921631011840$ gives
the number $n_G(p,q)$ in the source convention. The pairs remain
ordered; choosing $p\leq q$ does not introduce a further duality quotient.
\begin{center}\begin{tabular}{rrrrr}\toprule
$p$&$q$&$x_G/|\operatorname{Aut}G|$&$(x_G-e_G)/|\operatorname{Aut}G|$&$n_G$\\\midrule
5 & 7 & $\frac{649}{4}$ & $\frac{49}{4}$ & 150\\
5 & 14 & $\frac{1165}{4}$ & $\frac{5}{4}$ & 290\\
5 & 28 & $562$ & $2$ & 560\\
7 & 10 & $\frac{6121}{4}$ & $\frac{97}{4}$ & 1506\\
7 & 15 & $1298$ & $14$ & 1284\\
7 & 20 & $2918$ & $2$ & 2916\\
10 & 14 & $\frac{10109}{4}$ & $\frac{5}{4}$ & 2526\\
10 & 28 & $4994$ & $2$ & 4992\\
14 & 15 & $2330$ & $0$ & 2330\\
14 & 20 & $5026$ & $2$ & 5024\\
15 & 28 & $4496$ & $0$ & 4496\\
20 & 28 & $10297$ & $5$ & 10292\\
\bottomrule\end{tabular}\end{center}

Here the two intermediate columns are normalized pair counts; they
need not be integers, since automorphisms need not act freely on
nongenerating pairs.
Every division has zero remainder. All twelve values agree with the
fixed Table 1 of Connor--Leemans. The entries $(14,15)$ and $(15,28)$
already had equal lower and upper bounds; the other ten become exact
through Lemma~\ref{unique}. This adds no new principal refutation.

The complete subgroup roster and character tables are borrowed;
the contribution is the uniform subgroup argument making these
particular bounds exact. No group-order search, new classification,
Lean verification, or conjecture settlement is claimed. The remaining
L046 source obligations are unaffected.
\begin{thebibliography}{9}
\bibitem{CL} T. Connor and D. Leemans, \emph{Algorithmic enumeration
of regular maps}, Ars Math. Contemp. 10 (2016), 211--222.
\bibitem{roster} T. Connor and D. Leemans, \emph{The subgroup lattice
of $O'N$}, fixed 23-page roster (edition identified in the accompanying
source bindings).
\bibitem{Wilson} R. A. Wilson, \emph{The maximal subgroups of the
O'Nan group}, Journal of Algebra 97 (1985), 467--473.
\end{thebibliography}
\end{document}
